\section{Reference implementation}\label{sec:implementation-new}

The following dependency-free code implements the sixteen-rule transform and its causal decoder.
The wiring index uses the truth-table convention of Section~\ref{sec:wirings}; selecting
\code{k=5} gives the historical $\cxor$ instance.

\begin{lstlisting}
def wiring(k, s, x):
    """W_k(s,x), where k's bits are w(0,0),w(0,1),w(1,0),w(1,1)."""
    row = 2 * s + x
    return (k >> (3 - row)) & 1

def encode(bits, m=4, k=5, table=None, prefix=None):
    table = list(table) if table is not None else [0] * (1 << m)
    context = list(prefix) if prefix is not None else [0] * m
    residual = []
    for x in bits:
        a = int("".join(map(str, context)), 2) if m else 0
        s = table[a]
        residual.append(x ^ s)
        table[a] = wiring(k, s, x)
        if m:
            context = context[1:] + [x]
    return residual

def decode(residual, m=4, k=5, table=None, prefix=None):
    table = list(table) if table is not None else [0] * (1 << m)
    context = list(prefix) if prefix is not None else [0] * m
    bits = []
    for r in residual:
        a = int("".join(map(str, context)), 2) if m else 0
        s = table[a]
        x = r ^ s
        bits.append(x)
        table[a] = wiring(k, s, x)
        if m:
            context = context[1:] + [x]
    return bits
\end{lstlisting}

The historical \code{cta.py} uses a \code{bitarray} table and performs the $\W5$ assignment only
when the residual is one.  As shown in Section~\ref{sec:cta}, omitting the assignment when the
residual is zero leaves exactly the same state.  Its inverse reconstructs $x=r\oplus s[a]$, writes
$x$ to the cell, and shifts the recovered bit into the context, matching \code{decode(..., k=5)}
above.
